Cieľom tejto práce je prispieť k zefektívneniu procesu ukladania
dát do DNA vytvorením modela narušenia, ktorý popisuje chyby vznikajúce pri čítaní a zápise dát.
Pochopením a formalizáciou týchto narušení je možné zvýšiť spoľahlivosť a presnosť ukladania
a načítavania informácií, čo je kľúčové pre praktické využitie DNA ako úložného média v budúcnosti.

Model narušenia opisuje pravdepodobnostné správanie chýb, ktoré vznikajú počas syntézy,
sekvenovania a uchovávania DNA molekúl. Tento model zachytáva typy chýb ako substitúcie,
inzercie a delécie nukleotidov, čím poskytuje základ pre návrh efektívnych kódovacích
a opravných schém pri ukladaní dát do DNA.
